\section{Conclusion}
\label{sec:conclusion}


We introduce \name, a plug-and-play framework for token-efficient video RAG that jointly allocates a shared token budget across frames and subtitles. Its Token Scheduler optimizes query relevance, diminishing token value, and evidence diversity through iterative water-filling, while incremental allocation and attention refinement enable efficient reuse across conversational turns. At comparable compute, \name improves accuracy by up to 4.2 percentage points and reduces multi-turn prefill compute by 2.5$\times$ under comparable accuracy targets.
% This paper introduces \name, a plug-and-play framework for token-efficient video RAG. For each query, \name jointly decides which frames and subtitles to retain and how many tokens each receives under a shared token budget. Its Token Scheduler weights a saturating token value function by query relevance, discounts redundant content, and solves the resulting objective by iterative water-filling, while incremental allocation and attention-guided relevance carry the allocation across the turns of a conversation. At matched compute, \name improves accuracy over Top-K by up to 4.2\%, and incremental allocation reduces prefill compute by 2.5$\times$ in multi-turn conversations.


\paragraph{Limitations and future work.} \name relies on a simple token value function fit against data via maximum likelihood. While this reinforces how a purely analytical framework can produce strong results against state-of-the-art methods, learned approaches (e.g. via MLPs) may nonetheless offer superior performance by considering additional factors such as question type. Deeper integration with graph-based RAG, by factoring in commonly used graphical inference methods like personalized pagerank may further enrichen \name.

% \paragraph{Conclusion.} 